\documentclass[12pt,a4paper]{report}

\usepackage[spanish,es-tabla]{babel}
\usepackage[utf8]{inputenc}
\usepackage[T1]{fontenc}
\usepackage{lmodern}
\usepackage{geometry}
\usepackage{setspace}
\usepackage{graphicx}
\usepackage{booktabs}
\usepackage{longtable}
\usepackage{array}
\usepackage{hyperref}
\usepackage{csquotes}
\usepackage{enumitem}

\geometry{left=3cm,right=2.5cm,top=3cm,bottom=3cm}
\onehalfspacing

\hypersetup{
    colorlinks=true,
    linkcolor=black,
    citecolor=black,
    urlcolor=blue
}

\begin{document}

\title{Prototipo de software para an\'alisis asistido de resonancias magn\'eticas lumbares sagitales}
\author{Enzo Asplanatti \\ Francisco Fabrello}
\date{2026}

\maketitle
\tableofcontents
\clearpage

\chapter{Introducci\'on}

El an\'alisis de im\'agenes m\'edicas constituye una de las \'areas donde la ingenier\'ia inform\'atica ha adquirido un rol cada vez m\'as relevante. En particular, la aplicaci\'on de t\'ecnicas de procesamiento de im\'agenes e inteligencia artificial permite asistir tareas vinculadas con la detecci\'on, segmentaci\'on, medici\'on y visualizaci\'on de estructuras anat\'omicas. Estas herramientas no buscan reemplazar el criterio profesional, sino complementar el trabajo del especialista mediante resultados cuantitativos, trazables y revisables.

Dentro de este contexto, la resonancia magn\'etica de columna lumbar representa una modalidad de estudio ampliamente utilizada para observar estructuras anat\'omicas como cuerpos vertebrales, discos intervertebrales y canal espinal. La revisi\'on de este tipo de im\'agenes requiere experiencia profesional, interpretaci\'on anat\'omica y, en muchos casos, mediciones o comparaciones visuales entre niveles. Sin embargo, parte del an\'alisis estructural contin\'ua dependiendo de observaci\'on manual, delimitaci\'on visual de regiones de inter\'es y mediciones realizadas de forma no completamente automatizada.

Esta situaci\'on presenta una oportunidad de desarrollo desde la ingenier\'ia inform\'atica: construir una herramienta que permita asistir el an\'alisis de resonancias magn\'eticas lumbares sagitales mediante segmentaci\'on autom\'atica de estructuras relevantes, extracci\'on de mediciones geom\'etricas simples, visualizaci\'on superpuesta de resultados y generaci\'on de una salida estructurada editable. El prop\'osito no es emitir diagn\'osticos ni reemplazar al profesional de salud, sino ofrecer un soporte operativo que permita organizar informaci\'on anat\'omica y cuantitativa de manera revisable.

El presente Proyecto Final de Ingenier\'ia propone el dise\~no y desarrollo de un prototipo funcional de software orientado al an\'alisis asistido de resonancias magn\'eticas lumbares sagitales. El sistema se plantea como una herramienta acad\'emica y experimental que trabaja sobre im\'agenes de RM lumbar, segmenta estructuras anat\'omicas previamente definidas, obtiene mediciones cuantitativas acotadas y presenta los resultados en una interfaz que permita su revisi\'on por parte de un profesional.

El proyecto se apoya principalmente en datasets p\'ublicos anotados, en particular SPIDER, que contiene resonancias magn\'eticas lumbares sagitales con m\'ascaras de referencia para v\'ertebras, discos intervertebrales y canal espinal. Esta decisi\'on permite desarrollar y evaluar t\'ecnicamente el prototipo sin depender inicialmente de datos cl\'inicos privados, favoreciendo la reproducibilidad del trabajo y reduciendo riesgos asociados al tratamiento de informaci\'on sensible.

Desde el punto de vista del alcance, el prototipo se limita a una primera versi\'on funcional o MVP. Este MVP contempla la carga y normalizaci\'on de estudios, segmentaci\'on de estructuras anat\'omicas seleccionadas, extracci\'on de mediciones geom\'etricas simples, visualizaci\'on superpuesta de las m\'ascaras sobre la imagen original y generaci\'on de una salida estructurada editable. Quedan fuera del alcance el diagn\'ostico cl\'inico aut\'onomo, la recomendaci\'on terap\'eutica, la validaci\'on cl\'inica multic\'entrica, la integraci\'on hospitalaria completa y la certificaci\'on regulatoria.

La organizaci\'on del documento sigue una progresi\'on desde la definici\'on del problema hacia el desarrollo y validaci\'on del prototipo. En primer lugar, se presenta el planteamiento del problema, la justificaci\'on, los objetivos, el alcance y la descripci\'on general de la soluci\'on propuesta. Luego se desarrollar\'an el estado del arte y el marco te\'orico, que fundamentar\'an t\'ecnica y conceptualmente el proyecto. Posteriormente se incorporar\'an el user research, los requerimientos, la arquitectura, la metodolog\'ia de desarrollo, la implementaci\'on, las pruebas, los resultados y las conclusiones.

De esta manera, el trabajo busca integrar conocimientos de procesamiento de im\'agenes m\'edicas, aprendizaje autom\'atico, dise\~no de software, interacci\'on humano-computadora y validaci\'on t\'ecnica, dentro de un prototipo acotado y orientado a un problema concreto: asistir el an\'alisis estructural de resonancias magn\'eticas lumbares sagitales mediante segmentaci\'on, medici\'on y revisi\'on profesional.

\chapter{Planteamiento del problema}

La interpretaci\'on de resonancias magn\'eticas lumbares es una tarea especializada que requiere conocimiento anat\'omico, experiencia cl\'inica y capacidad para analizar distintas estructuras en m\'ultiples niveles de la columna. En el flujo habitual de revisi\'on, el profesional debe acceder al estudio, recorrer las secuencias disponibles, identificar estructuras relevantes, observar relaciones anat\'omicas, detectar posibles alteraciones y, cuando corresponde, realizar mediciones o comparaciones visuales.

Aunque los sistemas actuales de visualizaci\'on m\'edica permiten acceder a las im\'agenes y realizar mediciones manuales, muchas tareas vinculadas con la delimitaci\'on de estructuras anat\'omicas y la obtenci\'on de valores cuantitativos contin\'uan dependiendo del trabajo manual o semimanual. Esto puede generar carga operativa, variabilidad entre observadores y dificultad para obtener resultados estructurados, trazables y f\'acilmente reutilizables.

En el caso de las resonancias magn\'eticas lumbares sagitales, estructuras como cuerpos vertebrales, discos intervertebrales y canal espinal son fundamentales para el an\'alisis anat\'omico. Sin embargo, su delimitaci\'on manual puede consumir tiempo y estar sujeta a diferencias de criterio. Adem\'as, cuando las mediciones no se encuentran integradas en una salida estructurada, su reutilizaci\'on o revisi\'on posterior puede depender de registros dispersos, capturas de pantalla o informes redactados manualmente.

El problema central que aborda este proyecto puede formularse de la siguiente manera:

\begin{displayquote}
Existe una oportunidad de asistencia inform\'atica en el an\'alisis estructural de resonancias magn\'eticas lumbares sagitales, dado que la segmentaci\'on de estructuras anat\'omicas, la obtenci\'on de mediciones simples y la generaci\'on de salidas revisables a\'un pueden requerir tareas manuales, variables y poco integradas dentro de un mismo flujo funcional.
\end{displayquote}

Este problema no implica afirmar que el profesional carezca de herramientas para analizar estudios de RM lumbar, ni que la interpretaci\'on cl\'inica pueda automatizarse por completo. Por el contrario, el proyecto reconoce que el criterio profesional es indispensable. La dificultad se ubica en tareas operativas espec\'ificas que pueden ser asistidas por software: delimitar estructuras, calcular mediciones geom\'etricas, visualizar resultados de manera clara y organizar una salida editable que facilite la revisi\'on.

A partir de esta problem\'atica, se identifican tres dimensiones principales.

En primer lugar, una dimensi\'on operativa: el an\'alisis manual de estructuras anat\'omicas puede demandar tiempo, especialmente cuando se requiere revisar varios niveles, comparar regiones o generar documentaci\'on complementaria. La automatizaci\'on parcial de segmentaciones y mediciones podr\'ia reducir parte de esta carga.

En segundo lugar, una dimensi\'on de variabilidad: las mediciones manuales o la delimitaci\'on visual de estructuras pueden diferir entre usuarios o entre distintos momentos de revisi\'on. Un sistema asistivo no elimina la necesidad de validaci\'on profesional, pero puede ofrecer una base inicial homog\'enea y trazable sobre la cual el profesional intervenga.

En tercer lugar, una dimensi\'on de documentaci\'on: los resultados generados durante la revisi\'on de una imagen no siempre quedan estructurados de forma editable, reutilizable y vinculada con la regi\'on anat\'omica de origen. Una salida organizada por estructura, nivel y medici\'on puede facilitar la trazabilidad del an\'alisis.

Por lo tanto, el desaf\'io de ingenier\'ia consiste en dise\~nar un prototipo que integre en un mismo flujo funcional la carga de im\'agenes, la segmentaci\'on autom\'atica de estructuras lumbares, la derivaci\'on de mediciones geom\'etricas, la visualizaci\'on de resultados y la generaci\'on de una salida editable. Este sistema debe mantener siempre al profesional dentro del circuito de revisi\'on, evitando presentar los resultados como diagn\'osticos autom\'aticos.

La soluci\'on propuesta se diferencia de un sistema cl\'inico completo o de un producto regulado. Su finalidad es acad\'emica y experimental. Busca demostrar la factibilidad de integrar componentes conocidos --dataset p\'ublico, modelo de segmentaci\'on, m\'etricas, mediciones, interfaz y salida estructurada-- en un prototipo funcional acotado al an\'alisis asistido de RM lumbar sagital.

\chapter{Justificaci\'on}

El desarrollo del presente proyecto se justifica desde una perspectiva t\'ecnica, acad\'emica y funcional.

Desde el punto de vista t\'ecnico, la disponibilidad de datasets p\'ublicos anotados y de modelos de segmentaci\'on m\'edica permite abordar problemas que anteriormente requer\'ian acceso exclusivo a datos privados o infraestructura cl\'inica compleja. En particular, el dataset SPIDER ofrece una base adecuada para entrenar o evaluar modelos de segmentaci\'on sobre resonancias magn\'eticas lumbares sagitales, incluyendo estructuras alineadas con el alcance del prototipo: v\'ertebras, discos intervertebrales y canal espinal. Esto permite construir una prueba de concepto t\'ecnicamente viable y reproducible.

Adem\'as, existen arquitecturas y frameworks de aprendizaje profundo ampliamente utilizados en segmentaci\'on m\'edica, como U-Net, nnU-Net y herramientas especializadas en im\'agenes biom\'edicas. Estos antecedentes permiten enfocar el valor del proyecto no en crear necesariamente una arquitectura completamente nueva, sino en integrar un flujo completo: carga de estudios, preprocesamiento, segmentaci\'on, medici\'on, visualizaci\'on y salida editable.

Desde el punto de vista acad\'emico, el proyecto resulta pertinente para un Proyecto Final de Ingenier\'ia porque combina investigaci\'on aplicada, dise\~no de software, procesamiento de datos, inteligencia artificial, validaci\'on t\'ecnica y documentaci\'on de resultados. La propuesta no se limita a entrenar un modelo aislado, sino que busca construir un prototipo funcional con componentes articulados y verificables. Esto permite aplicar conocimientos propios de la carrera en un problema concreto y con restricciones reales.

A su vez, el proyecto se plantea dentro de l\'imites claros. No se propone desarrollar un sistema de diagn\'ostico aut\'onomo, ni una herramienta certificada para uso cl\'inico, ni una plataforma hospitalaria completa. Esta delimitaci\'on es importante porque evita sobreprometer capacidades y permite concentrar el esfuerzo en un MVP t\'ecnicamente alcanzable: segmentar estructuras, obtener mediciones simples y presentar resultados revisables.

Desde el punto de vista funcional, la herramienta propuesta puede aportar valor como soporte al an\'alisis estructural de RM lumbar. La segmentaci\'on autom\'atica permitir\'ia identificar visualmente regiones de inter\'es; las mediciones geom\'etricas aportar\'ian informaci\'on cuantitativa derivada de las m\'ascaras; y la salida estructurada editable permitir\'ia organizar los resultados para revisi\'on profesional. En todos los casos, el profesional conserva el control final sobre la interpretaci\'on, correcci\'on y eventual descarte de los resultados.

El estado del arte muestra que existen avances relevantes en segmentaci\'on autom\'atica de columna lumbar, datasets p\'ublicos y soluciones comerciales orientadas a mediciones y reportes. Sin embargo, muchos antecedentes se concentran en componentes parciales: algunos aportan datos, otros modelos de segmentaci\'on, otros m\'etricas de evaluaci\'on y otros productos cerrados. La brecha que justifica este PFI se encuentra en integrar esos elementos en un prototipo acad\'emico, acotado, reproducible y orientado a revisi\'on profesional.

El proyecto tambi\'en se justifica por su enfoque de trazabilidad. Cada medici\'on generada debe estar asociada a una estructura segmentada visible, y cada resultado debe poder ser revisado por el usuario. Esta caracter\'istica permite diferenciar el prototipo de una \enquote{caja negra} diagn\'ostica y orientarlo hacia una herramienta asistiva, donde la inteligencia artificial act\'ua como apoyo operativo.

Finalmente, la elecci\'on de trabajar inicialmente con datasets p\'ublicos reduce riesgos vinculados con privacidad y datos sensibles. Si en etapas futuras se incorporaran im\'agenes cl\'inicas no p\'ublicas, deber\'ian contemplarse procedimientos de anonimizaci\'on, consentimiento y resguardo normativo. Para el alcance actual, el uso de datos p\'ublicos permite concentrar el desarrollo en la validaci\'on t\'ecnica del pipeline y en la construcci\'on del prototipo funcional.

En s\'intesis, el proyecto se justifica porque aborda una necesidad real de asistencia en el an\'alisis estructural de RM lumbar, utiliza recursos t\'ecnicos disponibles, mantiene un alcance acad\'emico responsable y propone una integraci\'on funcional que conecta segmentaci\'on, medici\'on, visualizaci\'on y revisi\'on profesional.

\chapter{Objetivos}

\section{Objetivo general}

Dise\~nar y desarrollar un prototipo funcional de software que analice resonancias magn\'eticas lumbares sagitales para segmentar estructuras anat\'omicas relevantes, obtener un conjunto acotado de mediciones cuantitativas y generar una salida visual y estructurada editable que asista al profesional sin reemplazar su criterio cl\'inico.

\section{Objetivos espec\'ificos}

Para alcanzar el objetivo general, se definen los siguientes objetivos espec\'ificos:

\begin{enumerate}
    \item Relevar antecedentes acad\'emicos, datasets p\'ublicos, herramientas abiertas y soluciones comerciales relacionadas con el an\'alisis asistido de resonancias magn\'eticas lumbares.
    \item Identificar la brecha funcional que justifica el desarrollo de un prototipo acad\'emico orientado a segmentaci\'on, mediciones, visualizaci\'on y salida editable.
    \item Definir el alcance del MVP, especificando las estructuras anat\'omicas contempladas, las funcionalidades incluidas y las exclusiones cl\'inicas, t\'ecnicas y regulatorias.
    \item Analizar el flujo de trabajo y las necesidades de potenciales usuarios mediante actividades de user research orientadas a profesionales vinculados con resonancias magn\'eticas lumbares.
    \item Dise\~nar un pipeline de procesamiento que contemple carga, normalizaci\'on, preprocesamiento, segmentaci\'on, postprocesamiento, mediciones y visualizaci\'on de resultados.
    \item Implementar o adaptar un modelo de segmentaci\'on para identificar estructuras anat\'omicas en RM lumbar sagital, priorizando v\'ertebras, discos intervertebrales y canal espinal.
    \item Desarrollar un m\'odulo de mediciones cuantitativas simples derivadas de las m\'ascaras de segmentaci\'on generadas por el sistema.
    \item Dise\~nar una interfaz que permita visualizar la imagen original junto con las segmentaciones superpuestas y las mediciones asociadas.
    \item Generar una salida estructurada editable que organice los resultados del sistema sin constituir un informe cl\'inico autom\'atico.
    \item Evaluar t\'ecnicamente el desempe\~no del prototipo mediante m\'etricas de segmentaci\'on y comparaci\'on contra anotaciones de referencia del dataset utilizado.
    \item Complementar la validaci\'on t\'ecnica con una revisi\'on cualitativa profesional orientada a valorar plausibilidad anat\'omica, utilidad de las mediciones, claridad visual y legibilidad de la salida generada.
    \item Documentar las decisiones de dise\~no, limitaciones, resultados obtenidos y posibles l\'ineas de trabajo futuro.
\end{enumerate}

\chapter{Alcance}

El alcance del proyecto se define en funci\'on de un MVP acad\'emico orientado al an\'alisis asistido de resonancias magn\'eticas lumbares sagitales. El sistema propuesto no busca reemplazar el criterio profesional ni emitir diagn\'osticos, sino automatizar tareas operativas espec\'ificas vinculadas con segmentaci\'on, medici\'on, visualizaci\'on y documentaci\'on estructurada.

\section{Funcionalidades incluidas}

El proyecto incluye el dise\~no y desarrollo de un prototipo funcional con las siguientes capacidades:

\begin{itemize}
    \item Carga de estudios de RM lumbar sagital provenientes de datasets p\'ublicos o fuentes compatibles con el entorno de desarrollo.
    \item Normalizaci\'on y preprocesamiento de las im\'agenes para su uso dentro del pipeline.
    \item Segmentaci\'on autom\'atica de un conjunto acotado de estructuras anat\'omicas.
    \item Visualizaci\'on de las m\'ascaras generadas sobre la imagen original.
    \item Extracci\'on de mediciones geom\'etricas simples derivadas de las segmentaciones.
    \item Organizaci\'on de resultados por estructura anat\'omica y, cuando sea posible, por nivel lumbar.
    \item Generaci\'on de una salida estructurada editable para revisi\'on profesional.
    \item Evaluaci\'on t\'ecnica de las segmentaciones mediante m\'etricas cuantitativas.
    \item Revisi\'on cualitativa complementaria por parte de un profesional del \'area.
\end{itemize}

\section{Estructuras anat\'omicas contempladas}

El MVP se concentrar\'a en tres estructuras anat\'omicas principales:

\begin{itemize}
    \item cuerpos vertebrales;
    \item discos intervertebrales;
    \item canal espinal.
\end{itemize}

Estas estructuras se seleccionan porque se encuentran alineadas con el dataset p\'ublico definido como base del proyecto y porque permiten derivar mediciones geom\'etricas simples que pueden ser revisadas visualmente.

\section{Mediciones contempladas}

Las mediciones incluidas ser\'an acotadas y de naturaleza geom\'etrica. Podr\'an incluir, seg\'un la calidad de las m\'ascaras obtenidas y la informaci\'on espacial disponible:

\begin{itemize}
    \item \'areas proyectadas de estructuras segmentadas;
    \item dimensiones relativas de discos intervertebrales;
    \item medidas aproximadas asociadas al canal espinal en el plano sagital;
    \item relaciones o comparaciones simples entre niveles.
\end{itemize}

Estas mediciones no ser\'an interpretadas autom\'aticamente como diagn\'osticos. Por ejemplo, una reducci\'on de dimensi\'on o \'area no ser\'a clasificada por el sistema como patolog\'ia, sino presentada como dato cuantitativo revisable.

\section{Validaci\'on incluida}

La validaci\'on se desarrollar\'a en dos niveles.

En primer lugar, se realizar\'a una validaci\'on t\'ecnica cuantitativa, comparando las m\'ascaras generadas por el sistema contra anotaciones de referencia disponibles en el dataset utilizado. Para ello podr\'an emplearse m\'etricas como Dice, IoU, distancia de Hausdorff o m\'etricas de superficie, seg\'un corresponda.

En segundo lugar, se realizar\'a una validaci\'on cualitativa complementaria con un profesional vinculado al an\'alisis de im\'agenes m\'edicas. Esta revisi\'on se orientar\'a a evaluar la plausibilidad anat\'omica de las segmentaciones, la utilidad de las mediciones, la claridad de la visualizaci\'on y la legibilidad de la salida estructurada editable.

\section{Exclusiones del alcance}

Quedan fuera del alcance del proyecto:

\begin{itemize}
    \item diagn\'ostico cl\'inico autom\'atico;
    \item clasificaci\'on integral de patolog\'ias degenerativas;
    \item recomendaci\'on terap\'eutica;
    \item reemplazo del criterio m\'edico;
    \item uso cl\'inico real del sistema;
    \item validaci\'on cl\'inica multic\'entrica;
    \item integraci\'on completa con sistemas hospitalarios;
    \item integraci\'on PACS/RIS obligatoria;
    \item certificaci\'on regulatoria como software m\'edico;
    \item comparaci\'on longitudinal obligatoria entre estudios de un mismo paciente;
    \item entrenamiento sobre datos privados no anonimizados;
    \item desarrollo de un producto comercial final.
\end{itemize}

Estas exclusiones permiten mantener el proyecto dentro de un alcance acad\'emico realista y t\'ecnicamente verificable.

\section{Alcance del documento final}

El documento final abarcar\'a la fundamentaci\'on del problema, estado del arte, marco te\'orico, user research, requerimientos, arquitectura, metodolog\'ia de desarrollo, cronograma, implementaci\'on, validaci\'on, resultados, discusi\'on, conclusiones y trabajos futuros. La documentaci\'on buscar\'a dejar trazabilidad entre las decisiones de dise\~no y los resultados obtenidos durante el desarrollo del prototipo.

\chapter{Descripci\'on de la soluci\'on propuesta}

La soluci\'on propuesta consiste en un prototipo funcional de software para el an\'alisis asistido de resonancias magn\'eticas lumbares sagitales. El sistema recibir\'a como entrada una imagen o serie de RM lumbar, aplicar\'a un pipeline de procesamiento para segmentar estructuras anat\'omicas relevantes, obtendr\'a mediciones cuantitativas simples a partir de las m\'ascaras generadas y presentar\'a los resultados en una interfaz visual y una salida estructurada editable.

La soluci\'on se plantea como una herramienta de asistencia. Esto significa que el sistema no emitir\'a diagn\'osticos, no clasificar\'a patolog\'ias de manera aut\'onoma y no recomendar\'a tratamientos. Su funci\'on ser\'a asistir tareas operativas: delimitar estructuras, calcular valores geom\'etricos, visualizar resultados y organizar informaci\'on para revisi\'on.

\section{Flujo general del sistema}

El flujo general del prototipo puede describirse en las siguientes etapas:

\begin{enumerate}
    \item Carga del estudio: el usuario incorpora una RM lumbar sagital compatible con el sistema o selecciona un estudio previamente disponible dentro del entorno de prueba.
    \item Normalizaci\'on y preprocesamiento: la imagen se prepara para su an\'alisis mediante procedimientos como ajuste de formato, normalizaci\'on de intensidades, reorientaci\'on, remuestreo o recorte, seg\'un las necesidades del modelo utilizado.
    \item Segmentaci\'on anat\'omica: el sistema aplica un modelo de segmentaci\'on para identificar las estructuras contempladas en el MVP: v\'ertebras, discos intervertebrales y canal espinal.
    \item Postprocesamiento: las m\'ascaras generadas se procesan para mejorar su consistencia, separar componentes, organizar estructuras y preparar la informaci\'on para mediciones y visualizaci\'on.
    \item Extracci\'on de mediciones: a partir de las m\'ascaras se calculan mediciones geom\'etricas simples, tales como \'areas proyectadas, dimensiones relativas o relaciones entre estructuras.
    \item Visualizaci\'on superpuesta: la interfaz presenta la imagen original junto con las m\'ascaras segmentadas, permitiendo al usuario revisar visualmente los resultados.
    \item Revisi\'on profesional: el usuario puede aceptar, observar, corregir o descartar resultados. Esta etapa mantiene al profesional dentro del circuito de decisi\'on.
    \item Salida estructurada editable: el sistema genera una plantilla organizada con los resultados obtenidos, incluyendo estructuras, mediciones, unidades, observaciones y estado de revisi\'on.
\end{enumerate}

\section{M\'odulo de entrada y preprocesamiento}

El m\'odulo de entrada tendr\'a como finalidad cargar las im\'agenes que ser\'an analizadas por el prototipo. En una primera etapa, el sistema trabajar\'a con datos provenientes de datasets p\'ublicos anotados, priorizando aquellos que permitan evaluar segmentaciones contra m\'ascaras de referencia.

El preprocesamiento buscar\'a homogeneizar la informaci\'on de entrada para facilitar la aplicaci\'on del modelo de segmentaci\'on. Debido a que las im\'agenes de resonancia magn\'etica pueden presentar variabilidad de intensidad, resoluci\'on, orientaci\'on y formato, esta etapa resulta fundamental para obtener resultados consistentes.

Entre las tareas posibles se incluyen normalizaci\'on de intensidades, ajuste de dimensiones, conversi\'on de formato, remuestreo espacial y preparaci\'on de tensores de entrada para el modelo. Las transformaciones aplicadas deber\'an documentarse para asegurar trazabilidad y reproducibilidad.

\section{M\'odulo de segmentaci\'on}

El m\'odulo de segmentaci\'on ser\'a responsable de identificar autom\'aticamente las estructuras anat\'omicas definidas en el alcance. Para ello se implementar\'a o adaptar\'a un modelo de aprendizaje profundo orientado a segmentaci\'on de im\'agenes m\'edicas.

El sistema podr\'a utilizar como referencia arquitecturas ampliamente empleadas en este campo, tales como U-Net, nnU-Net u otros enfoques compatibles con el dataset y las restricciones de hardware disponibles. La elecci\'on final depender\'a de la factibilidad de entrenamiento, disponibilidad de recursos, rendimiento obtenido y facilidad de integraci\'on con el resto del prototipo.

La salida principal de este m\'odulo ser\'a una m\'ascara anat\'omica o conjunto de m\'ascaras que indiquen las regiones correspondientes a v\'ertebras, discos intervertebrales y canal espinal. Estas m\'ascaras ser\'an utilizadas posteriormente para visualizaci\'on, mediciones y evaluaci\'on t\'ecnica.

\section{M\'odulo de mediciones cuantitativas}

El m\'odulo de mediciones utilizar\'a las m\'ascaras generadas para obtener valores geom\'etricos simples. La finalidad de estas mediciones ser\'a aportar informaci\'on cuantitativa de apoyo, no establecer diagn\'osticos cl\'inicos.

Cada medici\'on deber\'a estar asociada a una estructura segmentada visible y, cuando corresponda, a un nivel anat\'omico. Adem\'as, deber\'a indicarse la unidad utilizada y conservarse la relaci\'on con la imagen y la m\'ascara de origen.

Las mediciones podr\'an incluir \'areas proyectadas, dimensiones relativas de discos, medidas asociadas al canal espinal en el plano sagital y comparaciones simples entre estructuras. La selecci\'on definitiva de mediciones depender\'a de la calidad de las m\'ascaras, la informaci\'on espacial disponible y la validaci\'on profesional.

\section{M\'odulo de visualizaci\'on}

El m\'odulo de visualizaci\'on permitir\'a observar la imagen original junto con las segmentaciones superpuestas. Esta funcionalidad resulta esencial porque permite revisar visualmente la plausibilidad anat\'omica de los resultados.

La interfaz deber\'a permitir diferenciar las estructuras segmentadas, activar o desactivar capas, consultar mediciones asociadas y revisar los resultados generados por el sistema. El dise\~no se orientar\'a a claridad y revisi\'on, no a reemplazo del visor m\'edico profesional.

La visualizaci\'on superpuesta tambi\'en servir\'a como soporte para la validaci\'on cualitativa, ya que el profesional podr\'a evaluar si los contornos generados resultan razonables, si las estructuras est\'an correctamente representadas y si la informaci\'on presentada es comprensible.

\section{M\'odulo de salida estructurada editable}

La salida estructurada editable ser\'a una plantilla generada a partir de los resultados del sistema. Esta salida podr\'a incluir datos como estructura anat\'omica, nivel, medici\'on, unidad, observaci\'on y estado de revisi\'on.

Su car\'acter editable es fundamental. El usuario debe poder corregir valores, agregar comentarios, descartar mediciones o indicar que un resultado requiere revisi\'on. De esta forma, la salida no se presenta como informe cl\'inico cerrado, sino como un documento de apoyo generado autom\'aticamente y supervisado por el profesional.

Este m\'odulo representa uno de los aportes principales del prototipo, porque conecta los resultados t\'ecnicos del modelo con una forma de documentaci\'on revisable y trazable.

\section{Rol del profesional}

El prototipo se dise\~nar\'a bajo un enfoque de humano en el circuito. El sistema propone resultados, pero el profesional conserva la decisi\'on final. Esta decisi\'on de dise\~no responde tanto a razones \'eticas como t\'ecnicas: las segmentaciones autom\'aticas pueden presentar errores, las mediciones requieren interpretaci\'on contextual y el alcance del proyecto no contempla diagn\'ostico aut\'onomo.

El profesional podr\'a intervenir revisando la segmentaci\'on, evaluando mediciones, modificando la salida editable y aportando observaciones. De esta manera, la inteligencia artificial funciona como soporte operativo y no como reemplazo del juicio experto.

\section{Resultado esperado del MVP}

El resultado esperado es un prototipo capaz de ejecutar un flujo completo sobre RM lumbar sagital: cargar una imagen, procesarla, segmentar estructuras anat\'omicas, generar mediciones, visualizar resultados y producir una salida editable.

El \'exito del MVP no depender\'a \'unicamente de alcanzar m\'etricas elevadas de segmentaci\'on, sino de demostrar una integraci\'on funcional coherente entre los m\'odulos. La validaci\'on deber\'a considerar tanto el rendimiento t\'ecnico como la utilidad percibida del flujo propuesto.

En consecuencia, la soluci\'on propuesta se ubica en un punto intermedio entre la investigaci\'on algor\'itmica y el desarrollo de software aplicado. No se limita a evaluar un modelo, pero tampoco pretende convertirse en un producto cl\'inico final. Su aporte consiste en construir una herramienta acad\'emica, trazable y revisable que permita explorar la asistencia inform\'atica en el an\'alisis estructural de resonancias magn\'eticas lumbares sagitales.

\section{Transici\'on hacia cap\'itulos 7 y 8}

Una vez definido el problema, los objetivos, el alcance y la soluci\'on propuesta, resulta necesario analizar qu\'e antecedentes existen en relaci\'on con el an\'alisis asistido de resonancias magn\'eticas lumbares, la segmentaci\'on autom\'atica de estructuras anat\'omicas, las mediciones cuantitativas derivadas de m\'ascaras y las herramientas de revisi\'on profesional. Por este motivo, el siguiente cap\'itulo desarrollar\'a el estado del arte, con el fin de identificar datasets, modelos, soluciones comerciales y brechas funcionales relevantes para el presente proyecto.

Posteriormente, el marco te\'orico establecer\'a los conceptos necesarios para comprender el dise\~no del prototipo: resonancia magn\'etica lumbar, anatom\'ia relevante, representaci\'on computacional de im\'agenes m\'edicas, segmentaci\'on, aprendizaje supervisado, m\'etricas de evaluaci\'on, visualizaci\'on, salida estructurada e inteligencia artificial asistiva.


% Capitulos a incorporar en una segunda instancia:
% \chapter{Estado del arte}

% Placeholder pendiente de migracion desde el documento existente.
% Este capitulo debera incorporar antecedentes academicos, datasets,
% modelos, soluciones comerciales y brechas funcionales vinculadas con el
% analisis asistido de resonancias magneticas lumbares sagitales.

% \chapter{Marco te\'orico}

% Placeholder pendiente de migracion desde el documento existente.
% Este capitulo debera desarrollar los conceptos necesarios para comprender
% el prototipo: RM lumbar, anatomia relevante, imagenes medicas,
% segmentacion, aprendizaje supervisado, metricas, visualizacion,
% salida estructurada e inteligencia artificial asistiva.


\nocite{*}
\bibliographystyle{plain}
\bibliography{references}

\end{document}
